%!TEX program = xelatex
% 如果你有参考文献（myref.bib），请按 xelatex -> bibtex -> xelatex*2 的顺序进行编译
\documentclass[dvipsnames, svgnames,a4paper,11pt]{article}
% Share-version redaction helpers
\newcommand{\answerblank}[1][3.8cm]{%
  \par\nopagebreak[4]\noindent\textbf{答：}\par\nopagebreak[4]\vspace*{#1}}
\newcommand{\recordblank}[1]{%
  \par\noindent\rule{0.95\linewidth}{0pt}\rule{0pt}{#1}\par}
\newcommand{\privateimage}[2]{%
  \rule{#1}{0pt}\rule{0pt}{#2}}
% ----------------------------------------------------
%   中山大学物理与天文学院本科实验报告模板
%   作者：Huanyu Shi，2019级
%   Github：https://github.com/huanyushi/SYSU-SPA-Labreport-Template
%   Last update : 2024.10.27
% ----------------------------------------------------

\input{settings} % 导入模板的相关设置
\usepackage{lipsum}


%---------------------------------------------------------------------
%	正文
%---------------------------------------------------------------------

\begin{document}

\scoresTable{20}{}{30}{}{30}{}{80}{}

% \infoTable{专业}{年级}{姓名}{学号}{实验时间}{教师签名}
\infoTable{\rule{2.5cm}{0.4pt}}{\rule{2.5cm}{0.4pt}}{\rule{2.5cm}{0.4pt}}{\rule{3cm}{0.4pt}}{\rule{3cm}{0.4pt}}{\rule{3cm}{0.4pt}}

\begin{center}
	\LARGE 大学物理实验 \quad TEC控温实验报告
\end{center}

\textbf{【实验报告注意事项】}
\begin{enumerate}
	\item 实验报告由三部分组成：
	\begin{enumerate}
		\item 预习报告：（提前一周）认真研读\underline{\textbf{实验讲义}}，弄清实验原理；实验所需的仪器设备、用具及其使用（强烈建议到实验室预习），完成课前预习思考题；了解实验需要测量的物理量，并根据要求提前准备实验记录表格（第一循环实验已由教师提供模板，可以打印）。预习成绩低于10分（共20分）者不能做实验。
	    \item 实验记录：认真、客观记录实验条件、实验过程中的现象以及数据。实验记录请用珠笔或者钢笔书写并签名（\textcolor{red}{\textbf{用铅笔记录的被认为无效}}）。\textcolor{red}{\textbf{保持原始记录，包括写错删除部分，如因误记需要修改记录，必须按规范修改。}}（不得输入电脑打印，但可扫描手记后打印扫描件）；离开前请实验教师检查记录并签名。
	    \item 分析讨论：处理实验原始数据（学习仪器使用类型的实验除外），对数据的可靠性和合理性进行分析；按规范呈现数据和结果（图、表），包括数据、图表按顺序编号及其引用；分析物理现象（含回答实验思考题，写出问题思考过程，必要时按规范引用数据）；最后得出结论。
	\end{enumerate}
	\textbf{实验报告就是将预习报告、实验记录、和数据处理与分析合起来，加上本页封面。}
	\item 每次完成实验后的一周内交\textbf{实验报告}（特殊情况不能超过两周）。
	\item 除实验记录外，实验报告其他部分建议双面打印。
\end{enumerate}


\clearpage


\setcounter{section}{0}
\nsection{大学物理实验-}{TEC控温实验}{预习报告}
	
\subsection{实验目的}
\begin{enumerate}
	\item 了解 TEC 半导体控温的原理和 PID 参数的调节。
	\item 基于 LabVIEW 和 myDAQ 设计和搭建温度采集和控制系统。
\end{enumerate}

\subsection{仪器用具}
\begin{table}[htbp]
	\centering
	\renewcommand\arraystretch{1.6}
	\begin{tabular}{p{0.05\textwidth}|p{0.20\textwidth}|p{0.05\textwidth}|p{0.5\textwidth}}
	\hline
	编号 & 仪器用具名称 & 数量 & 主要参数（型号，测量范围，测量精度等） \\
	\hline
	1 & 计算机       & 1   & Windows 系统，安装 LabVIEW 程序 \\
	2 & myDAQ        & 1   & NI myDAQ 数据采集与虚拟仪器平台 \\
	3 & 直流稳压电源 & 1   & DP832；TEC供电通道：电压 $12\,\mathrm{V}$，最大电流 $2\,\mathrm{A}$；变送器供电通道：$24\,\mathrm{V}$ \\
	4 & 控温实验平台 & 1   & 含防风罩、底座热沉、被控物体 \\
	5 & IBT-2 电机驱动器 & 1 & 输入电压 $6\sim27\,\mathrm{V}$，最大电流 $43\,\mathrm{A}$，PWM 控制，控制电平 $3.3\sim5\,\mathrm{V}$ \\
	6 & 温度变送器   & 1   & $24\,\mathrm{V}$ 供电，测温范围 $0\sim100\,^\circ\mathrm{C}$，输出电压 $0\sim10\,\mathrm{V}$ \\
	7 & PT100 热敏电阻 & 1 & 铂电阻，$R_0=100\,\Omega$（$0\,^\circ\mathrm{C}$），$A=0.385\,\Omega/^\circ\mathrm{C}$ \\
	8 & TEC 半导体制冷片 & 1 & 型号 TEC1-12703，最大电流 $3.5\,\mathrm{A}$，最大电压 $15.4\,\mathrm{V}$，最大温差 $68\,^\circ\mathrm{C}$ \\
	9 & 电路元器件   & 1套 & 面包板、杜邦线、粗导线、鳄鱼夹线缆等 \\
	\hline
\end{tabular}
\end{table}

\subsection{原理概述}

\subsubsection{PID 反馈控制原理}

PID（比例-积分-微分）控制是一种广泛应用于工业和实验室的闭环反馈控制方法。其基本思路是：将被控量（如温度）的\textbf{实际测量值}与\textbf{设定值}做比较，得到误差信号 $e(t)=T_\mathrm{set}-T_\mathrm{meas}$，再对该误差信号进行比例、积分、微分三种运算，将加权结果叠加后输出给执行机构（本实验中为 TEC 驱动电路），从而将被控量锁定在设定值附近。控制器输出 $u(t)$ 的连续表达式为：
\begin{equation}
	u(t) = K_p\!\left[e(t) + \frac{1}{T_i}\int_0^t e(\tau)\,\mathrm{d}\tau + T_d\frac{\mathrm{d}e(t)}{\mathrm{d}t}\right],
	\label{eq:PID}
\end{equation}
其中 $K_p$ 为比例增益，$T_i$ 为积分时间，$T_d$ 为微分时间。三个分量的作用分别如下：

\begin{itemize}
	\item \textbf{比例控制（P）}：输出 $P = K_p e(t)$ 与当前误差成正比。$K_p$ 越大，系统对误差响应越快，但过大时会引起系统在设定值附近持续振荡。
	\item \textbf{积分控制（I）}：输出 $I = \frac{K_p}{T_i}\int_0^t e(\tau)\,\mathrm{d}\tau$ 与历史累积误差成正比，能够消除纯比例控制下的稳态误差。$T_i$ 过小（积分作用过强）会导致严重超调甚至振荡；$T_i$ 过大则响应迟缓。
	\item \textbf{微分控制（D）}：输出 $D = K_p T_d\frac{\mathrm{d}e}{\mathrm{d}t}$ 反映误差的变化速率，能超前预测趋势、抑制超调和振荡。$T_d$ 过大会使系统对噪声极为敏感；$T_d$ 过小则抑制超调效果不明显。
\end{itemize}

在数字计算机（LabVIEW）中，PID 采用离散采样实现，其输出表达式为：
\begin{equation}
	\mathit{PID}_\mathrm{output} = K_c\!\left[e(t) + \frac{1}{T_i}\sum\frac{e(t)_{n-1}+e(t)_n}{2}\,\Delta t + T_d\frac{e(t)_{n-1}-e(t)_n}{\Delta t}\right],
\end{equation}
输出范围默认为 $(-100,\,100)$，对应 PWM 占空比。实验中 $T_i$、$T_d$ 单位为分钟（min），$\Delta t$ 单位为秒（s），计算时需注意单位换算。

\textbf{手动整定 PID 参数的一般步骤（Ziegler-Nichols 法）：}
\begin{enumerate}
	\item 令 $T_i\to\infty$、$T_d=0$，逐步增大 $K_c$，直到系统输出出现持续等幅振荡，记录此时的临界增益 $K_{cR}$ 和振荡周期 $T_u$。
	\item 根据控制类型，按下表选取参数初始值：

	\begin{table}[H]
		\centering
		\renewcommand\arraystretch{1.5}
		\begin{tabular}{|c|c|c|c|}
		\hline
		控制类型 & $K_p$ & $T_i$ & $T_d$ \\
		\hline
		P   & $0.5\,K_{cR}$  & —               & — \\
		PI  & $0.45\,K_{cR}$ & $0.85\,T_u$     & — \\
		PID & $0.6\,K_{cR}$  & $0.5\,T_u$      & $0.125\,T_u$ \\
		\hline
		\end{tabular}
		\caption{Ziegler-Nichols 方法 PID 参数整定参考表}
	\end{table}

	\item 在初始值基础上进一步微调，直到系统在最小超调量和振荡下快速稳定至设定值。
\end{enumerate}

\subsubsection{半导体制冷器（TEC）结构及工作原理}

\paragraph{Peltier 效应}
当直流电流 $I$ 流过由两种不同材料（A 和 B）构成的回路时，在两种材料的连接处会出现吸热或放热现象，吸/放热速率为：
\begin{equation}
	q = \pi_{AB}\cdot I,
\end{equation}
其中 $\pi_{AB}$ 为 Peltier 系数。改变电流方向即可改变热流方向，实现加热与制冷的切换。这一现象称为 \textbf{Peltier 效应}。

\paragraph{Seebeck 效应}
Seebeck 效应是 Peltier 效应的逆效应：两种材料组成的回路中，若两个连接处存在温差 $\Delta T$，则会在回路中产生温差电动势：
\begin{equation}
	V = \alpha_{AB}\,\Delta T,
\end{equation}
其中 $\alpha_{AB}$ 为 Seebeck 系数。Peltier 系数与 Seebeck 系数的关系为 $\pi_{AB}=\alpha_{AB}T$（$T$ 为绝对温度）。

\paragraph{TEC 结构与控温原理}
本实验使用的 TEC（Thermo Electric Cooler，热电制冷器）由多对 P 型与 N 型半导体温差电偶串联而成，上下两侧用陶瓷片封装。在 N 型材料中热流方向与电流方向相反，在 P 型材料中热流方向与电流方向相同；二者并联排列后，热量可在 TEC 两端面之间定向传输：一侧吸热（冷面），另一侧放热（热面）并由散热器散去。改变通入 TEC 的电流方向即可互换冷热面，实现加热或制冷。

本实验中 TEC 由 IBT-2 电机驱动器驱动，通过 myDAQ 输出的 PWM 信号控制驱动器的占空比，进而调节 TEC 的平均功率，实现闭环温度控制。PT100 铂电阻作为温度传感器，其阻值与温度的关系为：
\begin{equation}
	R_T = R_0 + A\,T, \quad R_0=100\,\Omega,\quad A=0.385\,\Omega/^\circ\mathrm{C}.
\end{equation}


\subsubsection{傅里叶变换与功率谱密度}

\textbf{傅里叶变换基本原理}：傅里叶变换是将时域信号变换到频域的重要数学工具。对于连续时间信号 $x(t)$，其傅里叶变换定义为：

\begin{equation}\label{eq:ft}
    X(f) = \int_{-\infty}^{\infty} x(t)\, e^{-j2\pi f t}\, dt
\end{equation}

其逆变换为：

\begin{equation}\label{eq:ift}
    x(t) = \int_{-\infty}^{\infty} X(f)\, e^{j2\pi f t}\, df
\end{equation}

对于离散采样数据（如本实验通过 myDAQ 采集的温度序列 $x[n]$，$n=0,1,\dots,N-1$），采用离散傅里叶变换（DFT）：

\begin{equation}\label{eq:dft}
    X[k] = \sum_{n=0}^{N-1} x[n]\, e^{-j2\pi k n / N}, \quad k = 0,1,\dots,N-1
\end{equation}

实际计算中通常使用快速傅里叶变换（FFT）算法实现。

傅里叶变换的物理意义在于：任何复杂的时域信号都可以分解为一系列不同频率、不同幅值和相位的正弦波的叠加。将温度信号的时域数据变换到频域后，可以分析温度波动中各频率成分的分布情况，从而评估控温系统的稳定性和噪声特性。

\textbf{功率谱密度（PSD）与幅度谱密度（ASD）}：

功率谱密度（Power Spectral Density, PSD）描述信号功率在频域上的分布，定义为：

\begin{equation}\label{eq:psd}
    P_{\!xx}(f) = \lim_{T \to \infty} \frac{1}{T} \left| X_T(f) \right|^2
\end{equation}

其中 $X_T(f)$ 是信号在有限时间段 $T$ 内的傅里叶变换。对于离散采样，PSD 可通过周期图法（Welch 方法）等估计。PSD 的单位为 $[\text{K}^2/\text{Hz}]$（温度信号的平方每赫兹）。

幅度谱密度（Amplitude Spectral Density, ASD）是 PSD 的平方根：

\begin{equation}\label{eq:asd}
    A(f) = \sqrt{P_{\!xx}(f)}
\end{equation}

ASD 的单位为 $[\text{K}/\sqrt{\text{Hz}}]$，其物理意义更加直观——它直接反映了在每个频率上温度波动的幅度大小。

在本实验中，通过分析控温数据的 PSD 或 ASD，可以：
\begin{enumerate}
    \item 识别温度波动的主要频率成分，判断系统是否存在周期性振荡；
    \item 评估控温系统的噪声水平，分析低频漂移和高频噪声的特性；
    \item 比较不同 PID 参数下控温效果的频域特性，为参数优化提供依据。
\end{enumerate}


\subsection{实验前思考题}
\begin{question}
	1. 热敏电阻测温点位置不同，是否会影响 PID 参数的设定，为什么？
\end{question}
\answerblank[3.8cm]

\begin{question}
	2. 被控物体材料不同，是否会影响 PID 参数的设定，为什么？
\end{question}
\answerblank[3.8cm]

\begin{question}
	3. 除了脉冲宽度调制控温，可以使用连续信号控温吗？如何实现？
\end{question}
\answerblank[3.8cm]

\clearpage
\begin{table}
	\renewcommand\arraystretch{1.7}
	\centering
	\begin{tabularx}{\textwidth}{|X|X|X|X|}
	\hline
	专业：& \rule{2.5cm}{0.4pt} &年级：& \rule{2.5cm}{0.4pt} \\
	\hline
	姓名： &\rule{2.5cm}{0.4pt} & 学号：&\rule{3cm}{0.4pt}\\
	\hline
	室温：& \rule{2.5cm}{0.4pt} & 实验地点： & \rule{2.5cm}{0.4pt} \\
	\hline
	学生签名：&\rule{2.5cm}{0.4pt} & 评分： &\\
	\hline
	实验时间：& \rule{3cm}{0.4pt} & 教师签名：& \rule{3cm}{0.4pt} \\
	\hline
	\end{tabularx}
\end{table}

\nsection{大学物理实验-}{TEC控温实验}{实验记录}

\subsection{实验内容和步骤}

\subsubsection{连接电路和测试}

\begin{enumerate}
	\item 利用万用表测量实验装置热敏电阻（PT100）的阻值，确认温度传感器工作正常后，利用线缆（一端为香蕉插头，一端为鳄鱼夹）连接到 myDAQ 的电阻测量端口（利用 myDAQ 的数字万用表功能）。

	\item 按讲义图 9 完成 IBT-2 电机驱动器和控温实验装置、myDAQ、直流稳压电源之间的连线：
	\begin{enumerate}
		\item TEC 的红色线接 IBT-2 的 M$+$，黑色线接 M$-$。
		\item 直流稳压电源（通道 1，电压 $12\,\mathrm{V}$，限流 $2\,\mathrm{A}$）通过粗导线接 IBT-2 电源端（注意正负极）；\textbf{程序运行后再打开电源输出}。
		\item 用杜邦线将 IBT-2 数字控制端与 myDAQ 按如下对应关系连接：
		\begin{itemize}
			\item RPWM（加热）$\longleftrightarrow$ myDAQ DIO 0
			\item LPWM（制冷）$\longleftrightarrow$ myDAQ DIO 1
			\item R\_EN 和 L\_EN 短接后接 PWM 信号（myDAQ AO 0）
			\item VCC $\longleftrightarrow$ myDAQ 5V
			\item GND $\longleftrightarrow$ myDAQ DGND
		\end{itemize}
		\item 温度变送器：$24\,\mathrm{V}$ 供电，PT100 接入变送器，输出接 myDAQ AI 0（$0\sim10\,\mathrm{V}$ 对应 $0\sim100\,^\circ\mathrm{C}$）。
	\end{enumerate}

	\item 在电脑桌面双击"NI MAX"图标，在"设备和接口"栏下确认 myDAQ 设备名称（如 myDAQ1），并保证与 LabVIEW 控温程序中调用的设备名称一致。
\end{enumerate}

% 若有实物接线照片，可在此插入
% \begin{figure}[H]
%     \centering
%     \includegraphics[width=0.7\linewidth]{figure/wiring_photo.jpg}
%     \caption{实验电路连接实物照片}
%     \label{fig:wiring}
% \end{figure}

\subsubsection{学习并利用提供的 LabVIEW 程序进行控温实验}

依次完成以下控温实验：

\begin{enumerate}
	\item \textbf{仅使用比例 P 控温：}令 $T_i\to\infty$、$T_d=0$，由小到大设定不同 $K_C$ 值，观察并记录控温结果，找出温度出现明显周期性振荡时对应的临界值 $K_{CR}$ 及振荡周期 $T_u$。

	\item \textbf{PI 控温实验：}利用 Ziegler-Nichols 法，取 $K_C=0.45\,K_{CR}$，$T_i=0.85\,T_u$，观察控温效果并记录。

	\item \textbf{PID 控温实验：}利用 Ziegler-Nichols 法，取 $K_C=0.6\,K_{CR}$，$T_i=0.5\,T_u$，$T_d=0.125\,T_u$，观察控温效果，在此基础上进一步微调参数，使温度波动小于 $\pm0.1\,^\circ\mathrm{C}$。
	

\end{enumerate}

\subsection{PID 参数调节记录}
\recordblank{5.5cm}
\subsection{实验电路连接与LabVIEW 编程截图}

\recordblank{6cm}

\clearpage

\subsection{控温结果截图}
\recordblank{6cm}


\subsection{实验报告记录}
\recordblank{5.5cm}
\clearpage

\begin{table}
	\renewcommand\arraystretch{1.7}
	\begin{tabularx}{\textwidth}{|X|X|X|X|}
	\hline
	专业：& \rule{2.5cm}{0.4pt} & 年级：& \rule{2.5cm}{0.4pt}\\
	\hline
	姓名：& \rule{2.5cm}{0.4pt} & 学号：& \rule{3cm}{0.4pt}\\
	\hline
	日期：& \rule{3cm}{0.4pt} & 评分：& \rule{3cm}{0.4pt}\\
	\hline
	\end{tabularx}
\end{table}

\nsection{大学物理实验-}{TEC控温}{分析与讨论}
\subsection{实验数据分析}
\recordblank{3.2cm}
\subsubsection{控温结果时域分析}
\recordblank{3.2cm}
\subsubsection{时域分析}
\recordblank{3.2cm}
\subsubsection{控温结果频域分析}
\recordblank{3.2cm}
\subsubsection{频域分析}
\recordblank{3.2cm}
\subsection{实验后思考题}

\begin{question}
	影响控温精度和稳定度的因素都有哪些，如何进行改进？
\end{question}
\answerblank[3.8cm]

\begin{question}
	本实验中直流电源的电流输出建议设置为 $2\,\mathrm{A}$，为了减小温度的波动，电流设置应增大还是减小，为什么？
\end{question}
\answerblank[3.8cm]

\begin{question}
	实验是否达到了控温的目的，还有哪些可以改进的地方？
\end{question}
\answerblank[3.8cm]

\begin{question}
	温度控制程序是否工作正常，是否还有优化的地方？
\end{question}
\answerblank[3.8cm]

\subsection{实验过程遇到问题、感受与想法}
\recordblank{5cm}
\clearpage
\appendix
\appendixpage
\addappheadtotoc
\subsection{附件（实验后收拾好桌面照片）}
\recordblank{5.5cm}
\end{document}
